\documentclass[11pt,a4paper]{article}
\usepackage[margin=2.5cm]{geometry}
\usepackage{amsmath,amssymb}
\usepackage{booktabs}
\usepackage{microtype}
\usepackage[hidelinks]{hyperref}

\newcommand{\W}{\mathsf{W}}
\newcommand{\Lb}{\mathsf{L}}
\newcommand{\D}{\mathsf{D}}

\title{How Wolf and Lamb was solved}
\author{}
\date{26 September 2026}

\begin{document}
\maketitle

\section*{The result}

Wolf and Lamb on the $5\times5$ board, with three wolves against fifteen lambs,
is a \textbf{draw} with best play. Every position of the game was computed and
independently checked. Only one of the eight opening moves keeps the draw: the
wolf capturing the central lamb, $(1,3)\times(3,3)$. Every other first move loses.
With 14 or 15 lambs on the board the wolves cannot win from any position, and
99.98\,\% of 15-lamb positions are lamb wins; the starting position is one of the
rare exceptions, because the wolves' opening double threats let them capture down
to a seven-lamb fortress that the lambs can never close. The draw does not depend
on the 100-turn rule: the same answer comes out with the rule removed.

\section*{The rules}

Cells are $(r,c)$, rows and columns from 1 to 5. Wolves start on
$(1,2),(1,3),(1,4)$; lambs fill rows 3 to 5. Wolves move first. A piece steps to an
orthogonally adjacent empty cell. A wolf may instead jump over one empty cell in a
straight line onto a lamb, which is removed; capturing is optional. A side with no
legal move loses. The wolves win when fewer than three lambs remain. It is a draw
after 100 plies without a capture, or on a threefold repetition.

\section*{How many positions}

A position is the set of wolf cells, the set of lamb cells, and the side to move.
Since lambs are only ever removed, positions with $k$ lambs (a \emph{layer}) lead
only to layers $k$ and $k-1$. The game was therefore solved one layer at a time,
from $k=3$ up to $k=15$; layers with fewer than three lambs are wolf wins by rule.

The rules are the same after any rotation or reflection of the board, so
positions that are mirror images share one table entry. The 2300 possible wolf
placements fall into 319 symmetry classes, and each layer stores
\[
N_k = 319 \cdot \binom{22}{k} \cdot 2
\]
entries (the 2 is the side to move), in total 2.6 billion entries, 5.2\,GB on disk,
each holding the value of the position and a distance.

\section*{The exact formula}

The 100-turn rule resets on every capture, so what matters in a layer is how fast
a side can force its way \emph{out} of the layer, or end the game, before the
clock runs out. For a side $X$ (wolves or lambs) and a position $s$ with $k$
lambs, let $T_X(s)$ be the least number of plies within which $X$ can force,
against any defence, either a terminal win or a capture into a $(k-1)$-lamb
position that is an $X$ win. For a successor $c$ of $s$ put
\[
t_X(c)=\begin{cases}
T_X(c) & \text{if $c$ is reached by a step (same layer)},\\
0 & \text{if $c$ is reached by a capture and } V_{k-1}(c)=X,\\
\infty & \text{if $c$ is reached by a capture and } V_{k-1}(c)\ne X,
\end{cases}
\]
where $V_{k-1}$ is the already known value of the layer below. Then $T_X$ is the
least solution of
\begin{equation}\label{eq:T}
T_X(s)=
\begin{cases}
0 & \text{no legal move and the opponent is to move ($X$ has won)},\\
\infty & \text{no legal move and $X$ is to move ($X$ has lost)},\\
1+\min_{c} t_X(c) & \text{$X$ to move},\\
1+\max_{c} t_X(c) & \text{opponent to move},
\end{cases}
\end{equation}
the minimum and maximum running over all successors $c$ of $s$.  In words: on
its own move $X$ picks the fastest way out, on the opponent's move the opponent
picks the slowest, and a capture is an immediate exit that is good for $X$ only
if the position it lands in is already known to be an $X$ win. ``Least solution''
means values are only finite when a finite forcing sequence exists, so cycles in
the position graph can never manufacture a win.

The value of a position with a fresh clock is then
\begin{equation}\label{eq:V}
V_k(s)=
\begin{cases}
\W & \text{if } T_{\W}(s)\le 100,\\
\Lb & \text{if } T_{\Lb}(s)\le 100,\\
\D & \text{otherwise.}
\end{cases}
\end{equation}
The two conditions cannot both hold: a strategy forcing a wolf exit and a strategy
forcing a lamb exit cannot both succeed in the same game.

\paragraph{Why this is the true value.}
If $T_X(s)\le 100$, side $X$ follows the moves that lower $T_X$ by one each turn;
the opponent cannot raise it by more than that, so an exit arrives within the
clock, the sequence never repeats a position (because $T_X$ keeps falling), and
after a capture the same argument continues in the layer below. If
$T_X(s)>100$, the opponent can hold $T_X$ up for 100 plies, at which point the
clock ends the game, or repetition ends it earlier, and either way $X$ has not
won. So $X$ wins with $c$ plies already on the clock exactly when
$T_X(s)\le 100-c$, and \eqref{eq:V} is the value with a fresh clock, which is the
situation at the start of the game and after every capture.

\section*{How it was computed}

For each layer and each side $X$, equation~\eqref{eq:T} was solved backwards by
breadth-first search, in C, on twelve threads. Positions with $T_X=0$ are found
first (the opponent is stuck), then those with $T_X=1$ (a winning capture is
available, or every move is such a capture), and so on: whenever a position gets
its value, its predecessors in the layer are updated. A predecessor with $X$ to
move takes value $d+1$ the first time one of its successors is valued $d$; a
predecessor with the opponent to move keeps a count of its unvalued successors and
takes value $d+1$ when the last one is valued. A position whose opponent has a
capture that is \emph{not} an $X$ win can never be forced and is marked as an
escape. Everything still unvalued after both passes is a draw. The whole
computation took about twelve minutes.

\section*{How it was checked}

The solver was not trusted. A second program, written separately in JavaScript
with its own indexing code and using the move generator that had been checked
move-for-move against a plain transcription of the rules on 300\,000 positions,
read every table entry and tested a local condition. Call a successor $c$ a
``$Z$-win within $n$'' if it is a capture into a $Z$-win of the layer below, or a
step to a position labelled $Z$ with distance at most $n$. For a position with
mover $M$ and opponent $Y$ the label must satisfy:

\begin{center}
\begin{tabular}{@{}ll@{}}
\toprule
label & required property \\
\midrule
no legal move & label is $Y$ with distance 0 \\
$M$ wins in $T$ & some successor is an $M$-win within $T-1$, none within $T-2$, $1\le T\le100$ \\
$Y$ wins in $T$ & all successors are $Y$-wins within $T-1$, some not within $T-2$, $1\le T\le100$ \\
draw & no successor is an $M$-win within 99, some successor is not a $Y$-win within 99 \\
\bottomrule
\end{tabular}
\end{center}

These conditions determine the labels completely. By induction on the distance,
every labelled win really is a forced win in exactly the stated number of plies
(soundness), and every position that really is a forced win within 100 plies must
carry that label with the right distance, because its winning successor is
already correctly labelled and the conditions leave no other option
(completeness). What remains is exactly the set of draws. The check passed on all
2,348,715,942 entries with zero errors, for both the normal tables and the
tables computed without the 100-turn rule. As a further guard, 3\,000 random
entries were compared with a brute-force search over the plain rules; all agreed.

\section*{What the tables say}

\begin{center}\small
\begin{tabular}{@{}rrrrrrr@{}}
\toprule
lambs & entries & wolf wins & lamb wins & draws & longest wolf win & longest lamb win \\
\midrule
3  & 982\,520     & 979\,706    & 7           & 2\,807       & 34 & 1 \\
5  & 16\,801\,092  & 16\,182\,453 & 2\,919       & 615\,720     & 88 & 23 \\
7  & 108\,807\,072 & 73\,329\,451 & 802\,636     & 34\,674\,985  & 70 & 99 \\
8  & 204\,013\,260 & 88\,101\,691 & 8\,814\,216   & 107\,097\,353 & 72 & 159 \\
9  & 317\,353\,960 & 69\,818\,174 & 53\,921\,205  & 193\,614\,581 & 60 & 113 \\
10 & 412\,560\,148 & 35\,109\,001 & 176\,467\,067 & 200\,984\,080 & 42 & 113 \\
11 & 450\,065\,616 & 9\,596\,170  & 322\,390\,318 & 118\,079\,128 & 24 & 55 \\
12 & 412\,560\,148 & 840\,621    & 371\,557\,736 & 40\,161\,791  & 7  & 43 \\
13 & 317\,353\,960 & 4\,164      & 309\,017\,226 & 8\,332\,570   & 1  & 48 \\
14 & 204\,013\,260 & 0          & 203\,094\,682 & 918\,578     & -- & 31 \\
15 & 108\,807\,072 & 0          & 108\,784\,334 & 22\,738      & -- & 24 \\
\bottomrule
\end{tabular}
\end{center}

``Longest win'' is the largest $T$ in the layer, in plies to the next capture or
the end of the game; lamb wins needing more than 100 plies in layers 8 to 10 are
draws under the rules. The wolves' wins are short tactical ones and vanish
entirely above 13 lambs; the lambs' wins are long squeezes.

\paragraph{The opening.} With the wolves to move at the start:

\begin{center}
\begin{tabular}{@{}ll@{}}
\toprule
wolf move & value \\
\midrule
$(1,3)\times(3,3)$ & \textbf{draw} \\
$(1,2)\times(3,2)$, $(1,4)\times(3,4)$ & lambs win, forced exit within 24 plies \\
$(1,2)\to(2,2)$, $(1,3)\to(2,3)$, $(1,4)\to(2,4)$ & lambs win within 20 plies \\
$(1,2)\to(1,1)$, $(1,4)\to(1,5)$ & lambs win within 18 plies \\
\bottomrule
\end{tabular}
\end{center}

After the central capture the lambs have four replies and all four draw, but two
of them, $(3,2)\to(2,2)$ and $(3,4)\to(2,4)$, leave the wolves a single drawing
answer out of nine. Best play then sees the wolves win several lambs through
double threats, since one wolf on the third rank and two on the first attack two
lambs at once and the lambs can block only one lane per move, until about seven
lambs remain and neither side can make progress.

\section*{Files}

\texttt{solver/solve.c} is the solver (\texttt{cc -O3 -o solver/solve
solver/solve.c -lpthread}, then \texttt{solver/solve 15 12});
\texttt{solver/verify.js} the independent checker; \texttt{solver/spotcheck.js}
the brute-force comparison; \texttt{solver/play.js} prints best play or plays
against the tables; \texttt{solver/server.js} feeds exact values to the web
board. The rules are transcribed in \texttt{engine.js} from \texttt{SPEC.md}.

\end{document}
